% Section 2.2, from lectures/L02/README.md: the objectives, the evaluation
% questions, and the next lecture.

\section{Review}
\label{c2:sec:review}

\needspace{6\baselineskip}
\subsection*{What you should be able to do now}
After this chapter, you should:

\begin{itemize}
  \item Understand how devicetree separates \emph{what hardware there is} from \emph{the code that uses it}.
  \item Be able to read a board’s devicetree, and the merged result of a build in \code{zephyr.dts}.
  \item Be able to add and change nodes and properties with an overlay.
  \item Be able to write a binding for a node of your own.
  \item Be able to get devices and properties from the devicetree in C, entirely at compile time.
  \item Be able to drive and read GPIO pins, and react to edges with a callback.
  \item Be able to trace a devicetree-related build error back to the node that caused it.
\end{itemize}

\needspace{6\baselineskip}
\subsection*{Questions to test yourself}
You should be able to explain:

\begin{enumerate}
  \item What is the difference between a node’s name, its label, and an alias?
  \item What does a binding add to a devicetree node, and what happens to a node whose binding is missing?
  \item Why must a \code{compatible} string’s vendor prefix be known to the build?
  \item How does an overlay change a board’s devicetree, and where can you see the merged result?
  \item Why is \code{DT_PROP(node, prop)} resolved at compile time, and what does that make impossible?
  \item What does \code{device_is_ready()} check, and why must a driver’s user call it?
  \item What does \code{GPIO_ACTIVE_LOW} in the devicetree change in the C code that uses the pin?
  \item What is the difference between \code{'__device_dts_ord_32' undeclared} and \code{undefined reference
    to '__device_dts_ord_32'}, and what does the number tell you?
\end{enumerate}

\needspace{6\baselineskip}
\subsection*{Looking ahead}
Kernel I, threads and time:

\begin{itemize}
  \item Threads, priorities, and cooperative and preemptive scheduling.
  \item Timers, and which context their expiry functions run in.
  \item Work queues, and deferring work out of interrupt context.
  \item What \code{native_sim}’s simulated time means for code that spins, and for stack sizes.
\end{itemize}

You will trace, line by line, which thread runs when, and why.
